\thispagestyle{fancy}
\fancyhead[R]{A.U 2025-2026}
\fancyhead[L]{Chapitre 1 : Cadre général du projet}

\chapter{Cadre général du projet}

\section*{Introduction}
Dans ce chapitre, nous commençons par la présentation du cadre de stage ainsi que
l'organisme d'accueil. Puis, nous mettons le projet dans son contexte général, en
identifiant la problématique et en étudiant les solutions existantes sur le marché.
Nous présentons ensuite la solution proposée. Finalement, nous exposons la
méthodologie de travail adoptée dans ce projet.

\section{Présentation de l'organisme d'accueil}
Fondée en 2018, CLEDISS~\cite{clediss} est une entreprise tunisienne éditrice de
logiciels, spécialisée dans la conception de solutions SaaS (\textit{Software as a
Service}) destinées à la digitalisation des processus métier. Créative et rigoureuse,
passionnée par le monde des technologies et des affaires, l'entreprise transforme les
ressources informatiques en véritables alliés stratégiques pour ses clients.

Positionnée sur le marché des progiciels de gestion intégrés, CLEDISS accompagne les
entreprises dans la transformation numérique de leurs activités en proposant une
plateforme hébergée dans le cloud, accessible à distance et maintenue en continu par
les équipes internes de l'entreprise.

L'équipe d'experts se consacre à l'élaboration de solutions sur mesure, adaptées aux
besoins spécifiques de chaque client. En travaillant en étroite collaboration avec eux,
l'entreprise comprend les défis uniques de chaque secteur et développe des stratégies
qui optimisent les opérations, augmentent la productivité et renforcent la
compétitivité sur le marché. La figure~\ref{fig:logo-clediss} présente le logo de
CLEDISS.

\begin{figure}[H]
    \centering
    \includegraphics[width=4cm]{Image/LogoClediss.png}
    \caption{Logo de l'entreprise CLEDISS}
    \label{fig:logo-clediss}
\end{figure}

\subsection{Fiche d'identité de CLEDISS}
Le tableau~\ref{tab:fiche-identite} présente la fiche d'identité de CLEDISS.

\begin{table}[H]
\centering
\renewcommand{\arraystretch}{1.4}
\begin{tabular}{|p{4.5cm}|p{9cm}|}
\hline
\textbf{Nom de la société} & \textbf{CLEDISS} \\
\hline
Année de création & 2018 \\
\hline
Secteur d'activité & Édition de logiciels SaaS \\
\hline
Produit phare & Nomadis ERP \\
\hline
Adresse & Tunis, Tunisie \\
\hline
Téléphone & +216 XX XXX XXX \\
\hline
Adresse-Mail & contact@clediss.com \\
\hline
Site Web & https://www.clediss.com \\
\hline
Logo & \includegraphics[width=3cm]{Image/LogoClediss.png} \\
\hline
\end{tabular}
\caption{Fiche d'identité de CLEDISS}~\cite{clediss}
\label{tab:fiche-identite}
\end{table}

\subsection{Organigramme de l'entreprise}
Dans cette partie, nous présentons l'organigramme de l'entreprise ainsi que le rôle de
chaque direction. La figure~\ref{fig:organigramme} illustre la structure
organisationnelle de CLEDISS.

\begin{figure}[H]
    \centering
    \includegraphics[width=0.9\textwidth]{Image/Organigramme.png}
    \caption{Organigramme de l'entreprise}
    \label{fig:organigramme}
\end{figure}

\subsection{Domaines d'activité de CLEDISS}
CLEDISS est une entreprise spécialisée dans l'édition de solutions logicielles SaaS et
les services informatiques destinés principalement aux entreprises souhaitant
digitaliser leurs processus métier. Ses domaines d'activité couvrent plusieurs axes
stratégiques :

\begin{itemize}[label=\ding{220}]
\item \textbf{Édition de progiciels de gestion intégrés :}
\begin{itemize}[label=\ding{43}]
    \item Conception et développement de la solution Nomadis ERP ;
    \item Gestion commerciale, suivi de stock, comptabilité et ressources humaines ;
    \item Évolutions fonctionnelles continues du produit.
\end{itemize}

\item \textbf{Digitalisation des processus métier :}
\begin{itemize}[label=\ding{43}]
    \item Automatisation de la force de vente et de la logistique ;
    \item Dématérialisation de la facturation et de l'inventaire ;
    \item Mise en place de solutions innovantes pour améliorer la productivité.
\end{itemize}

\item \textbf{Hébergement et exploitation d'infrastructure :}
\begin{itemize}[label=\ding{43}]
    \item Hébergement cloud de la plateforme Nomadis ERP ;
    \item Administration des serveurs et des bases de données ;
    \item Garantie de la continuité de service pour l'ensemble des clients.
\end{itemize}

\item \textbf{Accompagnement et support client :}
\begin{itemize}[label=\ding{43}]
    \item Paramétrage et accompagnement au déploiement de la solution ;
    \item Support technique et maintenance applicative ;
    \item Formation des utilisateurs sur les solutions déployées.
\end{itemize}
\end{itemize}

\section{Contexte du projet}
Dans cette section, nous décrivons la problématique, puis nous examinons les solutions
existantes ainsi que la solution proposée.

\subsection{Problématique}
Dans les entreprises éditrices de solutions SaaS, et plus particulièrement chez
CLEDISS, la disponibilité de l'infrastructure serveur constitue un élément central pour
assurer la continuité des services et la satisfaction des clients. Chaque incident
nécessite une détection rapide ainsi qu'une intervention efficace afin de garantir un
rétablissement dans les meilleurs délais.

Cependant, de nombreuses organisations continuent d'utiliser des méthodes
traditionnelles telles que la connexion manuelle en SSH, la consultation ponctuelle des
journaux système ou le suivi dans des fichiers Excel, voire l'absence totale de
supervision automatisée. Ces approches rendent la supervision peu centralisée, difficile
à suivre et souvent sujette à des retards ou à un manque de visibilité.

Cette absence de système unifié entraîne des difficultés dans la détection des
anomalies, la coordination entre les équipes techniques et le contrôle global du parc de
serveurs. Elle limite également la capacité à disposer d'indicateurs fiables pour
évaluer et améliorer la performance de l'infrastructure.

Les principales difficultés identifiées sont les suivantes :

\begin{itemize}[label=\ding{220}]
\item \textbf{Détection tardive des incidents :}
\begin{itemize}[label=\ding{55}]
    \item Une saturation de ressources (CPU, mémoire, disque) n'est identifiée qu'a
          posteriori, une fois ses effets visibles côté utilisateur ;
    \item Aucun mécanisme ne notifie automatiquement un administrateur en cas de
          dépassement d'un seuil critique.
\end{itemize}

\item \textbf{Absence de traçabilité :}
\begin{itemize}[label=\ding{55}]
    \item Aucune historisation des métriques permettant d'analyser l'évolution de la
          charge dans le temps ;
    \item Aucun journal des interventions effectuées sur les serveurs.
\end{itemize}

\item \textbf{Réactivité limitée :}
\begin{itemize}[label=\ding{55}]
    \item Toute action corrective nécessite une connexion manuelle, ce qui allonge le
          délai de résolution ;
    \item L'intervention exige un poste de travail équipé d'un client SSH, ce qui limite
          la réactivité en dehors des horaires de bureau.
\end{itemize}

\item \textbf{Absence de vision consolidée :}
\begin{itemize}[label=\ding{55}]
    \item Aucun tableau de bord unique ne permet de visualiser l'état de l'ensemble du
          parc de serveurs ;
    \item Aucune application mobile ne permet d'intervenir à distance en dehors du
          bureau.
\end{itemize}

\item \textbf{Absence de sauvegarde automatisée :}
\begin{itemize}[label=\ding{55}]
    \item Les sauvegardes, lorsqu'elles existent, ne sont ni planifiées, ni historisées,
          ni notifiées ;
    \item Aucun email de confirmation n'est envoyé après l'exécution d'une sauvegarde.
\end{itemize}
\end{itemize}

Ces insuffisances montrent la nécessité de mettre en place une solution moderne,
centralisée et automatisée, déployée selon une approche DevSecOps garantissant qualité,
sécurité et reproductibilité des mises en production.

La problématique peut donc être résumée comme suit :

\begin{center}
\fbox{\parbox{0.92\textwidth}{\centering \textbf{Comment surveiller en temps réel
l'état des serveurs, détecter les anomalies et intervenir à distance de manière
sécurisée ? Et comment automatiser et sécuriser le déploiement de cette solution, du
code jusqu'à la mise en production ?}}}
\end{center}

\subsection{Étude de l'existant}
Dans le domaine de la supervision d'infrastructure, plusieurs solutions de monitoring
permettent de collecter des métriques, de générer des alertes et de visualiser l'état
des serveurs. Ces outils visent à améliorer la disponibilité des services, la
traçabilité des opérations ainsi que la productivité des équipes techniques.

\subsubsection*{Zabbix}
Zabbix~\cite{zabbix} est une solution open source dédiée à la supervision
d'infrastructure. Elle permet la collecte de métriques, la génération d'alertes,
l'inventaire des équipements ainsi que la gestion des utilisateurs.

\subsubsection*{Nagios}
Nagios~\cite{nagios} est un outil de supervision réseau et système très répandu. Il
offre un système d'alertes par courrier électronique et une grande stabilité, reconnue
dans les environnements de production.

\subsubsection*{Prometheus et Grafana}
Prometheus~\cite{prometheus} associé à Grafana~\cite{grafana} constitue une combinaison
moderne de collecte de métriques et de visualisation. Elle est largement utilisée dans
les environnements conteneurisés et orchestrés par Kubernetes.

\subsubsection*{Comparaison des solutions existantes}
Le tableau~\ref{tab:comparaison-existant} présente une comparaison des principales
solutions de supervision étudiées.

\begin{table}[H]
\centering
\renewcommand{\arraystretch}{1.5}
\begin{tabular}{|p{4.8cm}|p{2.6cm}|p{2.6cm}|p{3.2cm}|}
\hline
\textbf{Critères} & \textbf{Zabbix} & \textbf{Nagios} & \textbf{Prometheus + Grafana} \\
\hline
Collecte de métriques & Oui & Oui & Oui \\
\hline
Alertes par email & Oui & Oui & Oui \\
\hline
Tableaux de bord & Moyen & Limité & Très avancé \\
\hline
Suivi en temps réel & Oui & Limité & Oui \\
\hline
Actions à distance (SSH) & Non & Non & Non \\
\hline
Application mobile native & Non & Non & Non \\
\hline
Gestion des sauvegardes & Non & Non & Non \\
\hline
Détection auto. des services & Partiel & Non & Partiel \\
\hline
Intégration DevSecOps & Non & Non & Partiel \\
\hline
Facilité de configuration & Faible & Faible & Moyenne \\
\hline
Modèle économique & Open Source & Open Source & Open Source \\
\hline
\end{tabular}
\caption{Comparaison des solutions existantes de supervision}
\label{tab:comparaison-existant}
\end{table}

\subsubsection*{Critique de l'existant}
L'étude des solutions existantes montre que les outils de supervision offrent des
fonctionnalités avancées pour la collecte et la visualisation des métriques. Cependant,
aucune d'entre elles ne répond de manière complète aux besoins spécifiques de CLEDISS :

\begin{itemize}[label=\ding{55}]
\item Aucune n'intègre d'actions à distance permettant de redémarrer un service ou un
      serveur directement depuis l'interface ;
\item Aucune ne propose une application mobile native permettant d'intervenir en
      mobilité ;
\item Aucune ne gère les sauvegardes automatiques avec suivi et historique ;
\item Aucune n'est intégrée dans une chaîne DevSecOps complète, de l'analyse du code
      jusqu'au déploiement.
\end{itemize}

Ces constats mettent en évidence la nécessité de disposer d'un outil de supervision
centralisé, temps réel et accessible depuis n'importe quel appareil, ce qui justifie le
développement d'une solution sur mesure répondant précisément aux exigences du projet.

\subsection{Solution proposée}
Afin de répondre aux problématiques identifiées dans la partie précédente, nous
proposons la conception et le développement d'une \textbf{plateforme web et mobile
centralisée de supervision et de gestion à distance d'infrastructure serveur}.

Cette solution vise à digitaliser et automatiser l'ensemble du processus de supervision,
depuis la collecte des métriques sur les serveurs jusqu'à l'intervention corrective,
tout en assurant un suivi en temps réel et une meilleure traçabilité des interventions.

Le principe fondamental de cette solution est qu'elle adopte une approche hybride
combinant la \textbf{supervision applicative} et l'\textbf{administration à distance},
le tout industrialisé selon une démarche \textbf{DevSecOps}.

\subsubsection*{Principales fonctionnalités de la solution}

\begin{itemize}[label=\ding{220}]
\item \textbf{Collecte automatisée des métriques :} un agent Python léger, installé sur
      chaque serveur supervisé, collecte les métriques système (CPU, RAM, disque,
      réseau, uptime) toutes les 5 secondes et les transmet au backend centralisé.

\item \textbf{Tableau de bord temps réel :} visualisation consolidée de l'ensemble des
      serveurs avec leur statut (OK, WARNING, CRITICAL, OFFLINE), les métriques
      globales et le nombre d'alertes actives, actualisée en continu via WebSocket.

\item \textbf{Système d'alertes automatiques :} génération d'alertes lors du dépassement
      de seuils configurables, accompagnée d'une notification par courrier électronique
      à l'administrateur.

\item \textbf{Actions à distance sécurisées :} redémarrage ou arrêt d'un service
      applicatif, redémarrage complet d'un serveur, exécutés via une connexion SSH
      initiée par le backend, sans ouverture manuelle de session.

\item \textbf{Détection automatique des services :} identification des services
      \texttt{systemd} réellement actifs sur chaque serveur, sans configuration manuelle
      préalable.

\item \textbf{Sauvegardes automatisées :} exécution planifiée quotidienne des
      sauvegardes, avec historique consultable et notification du résultat par email.

\item \textbf{Journalisation et audit :} traçabilité complète de toute action
      d'administration exécutée à distance (auteur, cible, commande, horodatage,
      résultat).

\item \textbf{Application mobile :} accès à l'ensemble des fonctionnalités depuis un
      smartphone, consommant le même backend que l'application web.
\end{itemize}

\subsubsection*{Apports de la solution}
La solution proposée permet de :

\begin{itemize}[label=\ding{51}]
\item Centraliser l'ensemble des informations liées aux serveurs sur une seule
      plateforme ;
\item Détecter les incidents de manière proactive grâce aux alertes automatiques ;
\item Réduire significativement le délai d'intervention en cas d'incident ;
\item Assurer une meilleure transparence et traçabilité des opérations ;
\item Offrir une mobilité totale grâce à l'application mobile ;
\item Garantir des déploiements fiables et sécurisés grâce à la chaîne DevSecOps.
\end{itemize}

Ainsi, la solution proposée constitue une réponse efficace aux limites des systèmes
existants et permet d'optimiser la gestion globale de l'infrastructure au sein de
l'entreprise.

\section{Étude des méthodologies de développement}
Le choix d'une méthodologie de développement constitue une étape essentielle dans la
réussite d'un projet logiciel. Une méthode adaptée permet d'améliorer l'organisation du
travail, de faciliter la collaboration entre les différents acteurs et d'assurer une
réalisation efficace des fonctionnalités attendues.

\subsection{Choix de la méthodologie}
Il existe plusieurs méthodologies de développement logiciel, chacune présentant des
avantages et des limites. Afin de sélectionner la plus adaptée à notre projet, nous
avons réalisé une étude comparative des principales méthodes, présentée dans le
tableau~\ref{tab:comparaison-methodo}.

\begin{table}[H]
\centering
\renewcommand{\arraystretch}{1.5}
\begin{tabular}{|p{2.2cm}|p{5.8cm}|p{5.8cm}|}
\hline
\textbf{Méthode} & \textbf{Forces} & \textbf{Faiblesses} \\
\hline
RUP &
\ding{51} Bonne traçabilité. \newline
\ding{51} Architecture bien structurée. &
\ding{55} Coût de mise en œuvre élevé. \newline
\ding{55} Complexité lors des premières phases. \\
\hline
XP &
\ding{51} Développement progressif. \newline
\ding{51} Tests continus. \newline
\ding{51} Coût relativement réduit. &
\ding{55} Forte implication de l'équipe requise. \newline
\ding{55} Difficile à appliquer sur de grands projets. \\
\hline
2TUP &
\ding{51} Adaptée à des projets de différentes tailles. \newline
\ding{51} Prise en compte des risques. \newline
\ding{51} Développement itératif. &
\ding{55} Principalement centrée sur le développement. \newline
\ding{55} Peu de modèles documentaires. \\
\hline
Scrum &
\ding{51} Forte implication du client. \newline
\ding{51} Travail collaboratif. \newline
\ding{51} Adaptation rapide aux changements. \newline
\ding{51} Livraison incrémentale. &
\ding{55} Gestion plus complexe avec de grandes équipes. \newline
\ding{55} Dépend fortement de l'implication des membres. \\
\hline
\end{tabular}
\caption{Étude comparative des méthodologies de développement}
\label{tab:comparaison-methodo}
\end{table}

\subsection{Méthodologie adoptée}
À l'issue de cette étude comparative, nous avons retenu la méthodologie agile
\textbf{Scrum}, qui répond le mieux aux besoins de notre projet.

Scrum repose sur un développement itératif et incrémental organisé en cycles courts
appelés \textit{sprints}. Cette approche favorise la collaboration entre les membres de
l'équipe et les parties prenantes, tout en permettant de s'adapter rapidement aux
évolutions des besoins. Elle offre également la possibilité de livrer régulièrement des
versions fonctionnelles de l'application afin de recueillir les retours des utilisateurs
et d'améliorer progressivement le produit.

Ce choix est d'autant plus pertinent que notre projet combine deux volets de nature
différente — le développement fonctionnel de la plateforme et la mise en place
progressive d'une chaîne d'exploitation DevSecOps — chacun pouvant être livré et validé
de manière incrémentale.

\begin{figure}[H]
    \centering
    \includegraphics[width=0.9\textwidth]{Image/scrum.png}
    \caption{Cycle de la méthodologie Scrum}
    \label{fig:scrum}
\end{figure}

\subsubsection*{Organisation de l'équipe Scrum}
La méthodologie Scrum s'appuie sur trois rôles principaux :

\noindent\textbf{Product Owner :} il représente les besoins des utilisateurs et définit
la vision du produit. Il est chargé de gérer et de prioriser le \textit{backlog} afin de
maximiser la valeur du projet. Ses principales responsabilités sont :

\begin{itemize}[label=\ding{43}]
\item Définir et prioriser les fonctionnalités ;
\item Planifier les différentes \textit{releases} ;
\item Ajuster les priorités en fonction des besoins ;
\item Valider les fonctionnalités développées.
\end{itemize}

\noindent\textbf{Scrum Master :} il veille au respect des principes Scrum et facilite la
collaboration entre les membres de l'équipe. Il accompagne également l'équipe dans la
résolution des obstacles pouvant ralentir le développement.

\noindent\textbf{Équipe de développement :} elle est responsable de la conception, de
l'implémentation et des tests des différentes fonctionnalités du système afin
d'atteindre les objectifs fixés pour chaque sprint.

\subsubsection*{Répartition des rôles}
Dans le cadre de notre projet, les rôles Scrum sont répartis comme indiqué dans le
tableau~\ref{tab:roles-scrum}.

\begin{table}[H]
\centering
\renewcommand{\arraystretch}{1.4}
\begin{tabular}{|p{5cm}|p{8cm}|}
\hline
\textbf{Rôle Scrum} & \textbf{Personne responsable} \\
\hline
Product Owner & Encadrant professionnel CLEDISS \\
\hline
Scrum Master & Encadrant académique ESPRIT \\
\hline
Équipe de développement & Mariem Chaabani \\
\hline
\end{tabular}
\caption{Répartition des rôles Scrum dans notre projet}
\label{tab:roles-scrum}
\end{table}

\section*{Conclusion}
Dans ce chapitre, nous avons présenté l'organisme d'accueil ainsi que le contexte
général du projet. Nous avons ensuite exposé la problématique, étudié les solutions
existantes et comparé plusieurs méthodologies de développement afin de justifier le
choix de Scrum. Le chapitre suivant sera consacré à la phase préparatoire du projet,
communément appelée Sprint~0.
